\documentclass[11pt,a4paper]{article}
\usepackage[T1]{fontenc}
\usepackage{lmodern,amsmath,amssymb,amsthm,mathtools,booktabs,microtype}
\usepackage[margin=27mm]{geometry}
\usepackage[colorlinks=true,linkcolor=blue,citecolor=blue,urlcolor=blue]{hyperref}
\newtheorem{theorem}{Theorem}[section]
\newtheorem{proposition}[theorem]{Proposition}
\newtheorem{lemma}[theorem]{Lemma}
\newtheorem{corollary}[theorem]{Corollary}
\newcommand{\Q}{\mathbb Q}\newcommand{\Z}{\mathbb Z}
\DeclareMathOperator{\Tr}{Tr}\DeclareMathOperator{\rank}{rank}
\setlength{\emergencystretch}{2em}
\title{Complex Multiplication and Cancellation in Elliptic Trace Heights}
\author{}\date{10 November 2026}
\begin{document}\maketitle
\begin{abstract}
For the curve $E_D:y^2=x^3-Dx$ over a number field $K$ not
containing $i$, adjoining $i$ doubles the Mordell--Weil rank.
If $P\in E_D(K)$ is non-torsion, the point $C=(1+i)P$ generates
an infinite cyclic subgroup whose nonidentity abscissas all have
relative trace zero, while their heights grow quadratically.
We also give four pairwise distinct doubled points, with distinct
abscissas, whose abscissa sum is zero and whose heights are
unbounded. The curve with $D=11-6\sqrt2$ provides an explicit
example over $\Q(\sqrt2,i)$; a two-prime reduction certificate
proves that the displayed point has infinite order. These results
obstruct an unrestricted transfer of the particular elliptic rank
and height inputs to a complex field. They do not give a
counterexample to Diophantine definability of a ring of integers
in its fraction field.
\end{abstract}

\section{The curve and its endomorphisms}

Let $K$ be a number field with $i\notin K$, let $D\in K^*$, and put
\[
 E_D:y^2=x^3-Dx,\qquad L=K(i).
\]
In particular, $K$ may be any number field with a real embedding.
The curve is nonsingular, with discriminant $64D^3$.
Let $\sigma$ be the nonidentity element of $\operatorname{Gal}(L/K)$.
Write $O$ for the identity of $E_D$. The automorphism
\begin{equation}\label{eq:iota}
             \iota(O)=O,\qquad \iota(x,y)=(-x,iy)
\end{equation}
is defined over $L$. It preserves the curve equation and the
chord-and-tangent law, and satisfies
\begin{equation}\label{eq:endomorphisms}
 \iota^2=[-1],\qquad \sigma(\iota Q)=-\iota(\sigma Q)
                         \quad(Q\in E_D(L)).
\end{equation}
The notation $(1+\iota)P$ means $P+\iota P$, and similarly for
$1-\iota$. Thus
\begin{equation}\label{eq:normendo}
                    (1-\iota)(1+\iota)=[2].
\end{equation}
These identities are identities of elliptic endomorphisms, not
statements about a chosen generator of a Mordell--Weil group.

\begin{lemma}\label{lem:coordinates}
If $P=(x,y)$ and $x\ne0$, then
\begin{equation}\label{eq:Ccoordinates}
 \begin{split}
 x((1+\iota)P)&=-\frac{i(x^2-D)}{2x},\\
 y((1+\iota)P)&=-\frac{(1+i)y(x^2+D)}{4x^2}.
 \end{split}
\end{equation}
The endomorphism $1+\iota$ is a degree-two isogeny with kernel
$\{O,(0,0)\}$.
\end{lemma}
\begin{proof}
The line through $(x,y)$ and $(-x,iy)$ has slope
$(1-i)y/(2x)$. Its third-intersection formula gives
\[
 x((1+\iota)P)
       =\left(\frac{(1-i)y}{2x}\right)^2
       =-\frac{i(x^2-D)}{2x}.
\]
Substituting into the ordinate formula gives the second identity.
If $(1+\iota)Q=O$, then $\iota Q=-Q$. For an affine point this
implies $-x(Q)=x(Q)$, hence $x(Q)=0$ and $Q=(0,0)$.
Conversely $(0,0)$ lies in the kernel. In characteristic zero a
nonzero elliptic endomorphism is separable, so its degree is the
number of geometric kernel points, namely two
\cite[Chapter III, Sections 4 and 6]{Silverman}.
\end{proof}

\section{Exact rank growth}

\begin{theorem}\label{thm:rank}
For every $K,D,L$ as above,
\[
                         \rank E_D(L)=2\rank E_D(K).
\]
More precisely,
\begin{equation}\label{eq:rankinclusion}
           2E_D(L)\subseteq E_D(K)+\iota E_D(K),
\end{equation}
and $E_D(K)\cap\iota E_D(K)$ consists of points killed by two.
\end{theorem}
\begin{proof}
For $Q\in E_D(L)$, put
\[
 A=Q+\sigma Q,\qquad B=-\iota(Q-\sigma Q).
\]
Both points are fixed by $\sigma$: the assertion for $B$ uses
\eqref{eq:endomorphisms}. Thus $A,B\in E_D(K)$, and
\[
                              2Q=A+\iota B.
\]
This proves~\eqref{eq:rankinclusion}.
If $T\in E_D(K)\cap\iota E_D(K)$, write $T=\iota U$ with
$U\in E_D(K)$. Then $T=\sigma T=-\iota U=-T$, so $2T=O$.
The sum $E_D(K)+\iota E_D(K)$ therefore has rank
$2\rank E_D(K)$. It has finite index in $E_D(L)$, since it
contains $2E_D(L)$ and the Mordell--Weil group is finitely generated.
The ranks are equal, proving the formula.
\end{proof}

In particular, if $E_D(K)$ has positive rank, the same curve
cannot retain that rank after adjoining $i$. If its rank over $K$
is one, its rank over $K(i)$ is two. Selecting a different elliptic
curve is a separate possibility: Zywina's theorem provides curves
of rank one stable in every prescribed quadratic extension
\cite[Theorem 1.2]{Zywina}. The obstruction here concerns the
specified family with $j=1728$.

\section{An infinite cyclic subgroup with zero trace}

Let $h$ denote the absolute logarithmic Weil height on
$\mathbb P^1(\overline{\Q})$, with $h(0)=h(\infty)=0$, and put
$x(O)=\infty$.
Normalize the canonical height on $E_D$ by
\[
 \widehat h(Q)=\frac12\lim_{m\to\infty}4^{-m}
                                      h(x([2^m]Q)).
\]
We use its quadraticity, Galois invariance, parallelogram law, and
positivity away from torsion, together with
\begin{equation}\label{eq:heightcomparison}
                  h(x(Q))=2\widehat h(Q)+O_{E_D}(1).
\end{equation}
These are the standard canonical-height properties
\cite[Chapter VIII, Section 9]{Silverman}.

\begin{theorem}\label{thm:trace}
Let $P\in E_D(K)$ have infinite order and put $C=(1+\iota)P$.
Then $C$ has infinite order. For every nonzero integer $n$,
\begin{equation}\label{eq:zerotrace}
 \Tr_{L/K}(x([n]C))=0,
 \qquad \Tr_{L/\Q}(x([n]C))=0,
\end{equation}
whereas
\begin{equation}\label{eq:tracegrowth}
                 h(x([n]C))=4n^2\widehat h(P)+O_{E_D}(1).
\end{equation}
In particular, even on the cyclic subgroup $\Z C$, the height of
an abscissa cannot be bounded by any function of the height of
its relative or absolute trace that is finite at zero.
The same conclusion holds on $2\Z C\subseteq2E_D(L)$.
\end{theorem}
\begin{proof}
Equation~\eqref{eq:normendo} sends $C$ to $[2]P$, so $C$ is
non-torsion. Since $\sigma P=P$,
\[
 \sigma C=(1-\iota)P=-\iota(1+\iota)P=-\iota C.
\]
For $Q=[n]C$ this gives
\[
 \sigma(x(Q))=x(\sigma Q)=x(-\iota Q)=-x(Q).
\]
The points are non-torsion, so their abscissas are finite.
This proves the relative trace identity; transitivity of trace
proves the absolute one.

The equality $x(\iota Q)=-x(Q)$ and the height definition give
$\widehat h(\iota Q)=\widehat h(Q)$. Galois invariance gives
$\widehat h(P+\iota P)=\widehat h(P-\iota P)$.
The parallelogram law therefore yields
\[
 2\widehat h(C)
       =\widehat h(P+\iota P)+\widehat h(P-\iota P)
       =2\widehat h(P)+2\widehat h(\iota P)
       =4\widehat h(P).
\]
Quadraticity and~\eqref{eq:heightcomparison} prove
\eqref{eq:tracegrowth}. Its leading coefficient is positive,
while both trace heights are identically zero. Taking only even
$n$ proves the last assertion.
\end{proof}

\section{Four distinct doubled points with cancelling abscissas}

\begin{theorem}\label{thm:sum}
For $n\ge1$, set
\[
 Q_{1,n}=[2n]P,\quad Q_{2,n}=\iota[2n]P,\quad
 Q_{3,n}=[4n]P,\quad Q_{4,n}=\iota[4n]P.
\]
These are four pairwise distinct non-torsion points in $2E_D(L)$.
Their abscissas are pairwise distinct and satisfy
\begin{equation}\label{eq:sumzero}
                       \sum_{j=1}^4x(Q_{j,n})=0.
\end{equation}
Nevertheless
\[
 \min_{1\le j\le4}h(x(Q_{j,n}))\longrightarrow\infty
                         \quad\text{as }n\longrightarrow\infty.
\]
Consequently there are no constants $A,B$ for which
\[
 \max_jh(x(Q_j))\le A\,h\!\left(\sum_jx(Q_j)\right)+B
\]
holds for all four-tuples of distinct non-torsion doubled points
of $E_D(L)$, even with pairwise distinct abscissas.
\end{theorem}
\begin{proof}
All four points are doubles and non-torsion. Their abscissas cancel
in pairs by~\eqref{eq:iota}.
If $aP=\iota bP$ for nonzero integers $a,b$, applying
$[a]+[b]\iota$ gives $[a^2+b^2]P=O$, a contradiction.
The same argument permits either sign on the right. Equal
abscissas on an elliptic curve mean equal or opposite points;
thus no abscissa from one of the two $\iota$-directions can equal
one from the other. Within either direction, equality would imply
$[2n]P=\pm[4n]P$, again impossible. This proves both distinctness
claims.
Finally~\eqref{eq:heightcomparison} gives
\[
 h(x([2n]P))=8n^2\widehat h(P)+O_{E_D}(1),\qquad
 h(x([4n]P))=32n^2\widehat h(P)+O_{E_D}(1).
\]
The $\iota$-images have the same heights, proving the assertions.
\end{proof}

\section{An explicit instance}

Take
\[
 K=\Q(\sqrt2),\qquad D=11-6\sqrt2,
\]
\[
 R_0=(1+\sqrt2,2\sqrt2),\qquad
 P=\left(\frac{57-28\sqrt2}{8},
          \frac{-660+417\sqrt2}{32}\right)=[2]R_0.
\]
Both curve identities and the doubling identity are exact.
The following short verification of non-torsion is independent
of any rank or Diophantine-definability input.

\begin{lemma}\label{lem:nontorsion}
The displayed point $P$ has infinite order.
\end{lemma}
\begin{proof}
Use the degree-one primes of $K$ specified by
$\sqrt2\equiv6\pmod{17}$ and $\sqrt2\equiv5\pmod{23}$.
Both are good-reduction primes. Direct finite-field calculation gives
\[
\begin{array}{c|c|c|c|c}
 p&D\bmod p&P\bmod p&\#E_D(\mathbb F_p)&
                   \operatorname{ord}(P\bmod p)\\\hline
 17&9&(1,14)&20&5\\
 23&4&(4,5)&24&3.
\end{array}
\]
If $P$ were torsion, its reduction of order five at the first
prime would force its order to be divisible by five. A point of
order five would then belong to $E_D(K)$. Reduction at the second
prime is injective on five-primary torsion
\cite[Chapter VII, Section 3]{Silverman}, but its group has order
$24$. This is impossible.
\end{proof}

Theorems~\ref{thm:rank}, \ref{thm:trace} and~\ref{thm:sum} now
apply over $\Q(\sqrt2,i)$ without additional arithmetic hypotheses.
The point $C=(1+\iota)P$ is itself a double because $P=[2]R_0$.
Thus all the trace-cancelling multiples in this explicit instance
already lie in $2E_D(\Q(\sqrt2,i))$.

\section{Scope for definitions over number fields}

A ring of integers $\mathcal O_L$ is Diophantine in $L$ if it
is the projection of a polynomial zero set with all witnesses
ranging over $L$. None of the preceding theorems disproves the
existence of such a definition. They rule out transferring the
particular rank-one and elliptic trace-height assertions to
unrestricted points of the larger field.

The cancellation in Theorem~\ref{thm:trace} is an explicit Galois
relation among the points. It violates the independence hypothesis
required by the aggregate trace-height theorem~\cite{Aggregate}. The cancellation
in Theorem~\ref{thm:sum} shows why a sum-height argument relying
on positive abscissas at real embeddings requires a different
justification over a field containing $i$.
The rank formula shows that even the condition of being a double
does not retain the original cyclic group.

The isogeny-multiplicity argument for modular contact heights~\cite{Multiplicity} is
separate. Its height averaging and dimension-two isogeny estimates
are formulated over arbitrary number fields; they do not require
real embeddings. The obstruction identified here is in the
elliptic group and trace or sum encoding. A contact-height theorem
for algebraic parameters would still require the appropriate
modular realization and its stated uniform estimates, rather than
an unrestricted change of its rational parameter field.

\section{Exact verification}

The standard-library program
\nolinkurl{computation/number_fields/check_trace_cancellation.py}
checks the coordinate identities in the universal ring
\[
 \Q[D,x,x^{-1},y,i]/(y^2-x^3+Dx,\ i^2+1).
\]
It verifies the automorphism, chord formulas, both coordinates of
$(1-\iota)(1+\iota)P=[2]P$ after clearing denominators, conjugation,
and the group decomposition used for rank growth. It also verifies
the displayed point and its doubling, and counts the two
finite reduction groups by two independent elementary methods.
Their point orders give the non-torsion certificate above.

Several exact multiples are included as supplementary checks of
the trace and four-point identities. The general identities and
height-growth conclusions follow from the proofs, not from those
finite samples. The output is
\nolinkurl{computation/number_fields/trace-cancellation-certificate.json}.

\begin{thebibliography}{9}
\bibitem{Aggregate}
\emph{Aggregate trace heights of elliptic multiples},
companion manuscript, 8 November 2026.
\nolinkurl{papers/trace_heights/aggregate_trace_height_20261108_v9.tex}.
\bibitem{Multiplicity}
\emph{Multiplicity of modular factors and quadratic contact-height bounds},
companion manuscript, 10 November 2026.
\nolinkurl{papers/contact_heights/multiplicity_and_contact_heights_20261110_v19.tex}.
\bibitem{Silverman}
J.~H.~Silverman, \emph{The Arithmetic of Elliptic Curves},
2nd edition, Graduate Texts in Mathematics~106, Springer, 2009;
Chapters III, VII and VIII.
\url{https://doi.org/10.1007/978-0-387-09494-6}.
\bibitem{Zywina}
D.~Zywina, \emph{Rank one elliptic curves and rank stability},
arXiv:2505.16960, Theorems~1.1 and~1.2.
\url{https://davidzywina.github.io/papers/Rank1WithStability.pdf}.
\end{thebibliography}
\end{document}
